% Figure: graphical determination of the critical Mach number. Two
% curves are drawn against freestream Mach number: the
% compressibility-corrected suction-peak pressure coefficient (rising in
% magnitude as M grows, Prandtl-Glauert), and the critical pressure
% coefficient Cp* (the Cp at which the local Mach number is exactly 1).
% Their intersection is M_cr. Cp is plotted as -Cp so both curves are
% positive and the crossing reads clean.
% Prandtl-Glauert: Cp = Cp0 / sqrt(1-M^2), here Cp0 = -0.43 (a thin
% section), so -Cp = 0.43/sqrt(1-M^2).
% Cp*: Cp* = (2/(g M^2)) [ ((1 + 0.2 M^2)/1.2)^3.5 - 1 ], g=1.4;
% -Cp* is the negative of a negative number, i.e. |Cp*|, plotted directly.
\begin{figure}[htbp]
\centering
\begin{tikzpicture}
\begin{axis}[
  width=11cm, height=8cm,
  xlabel={freestream Mach number $M_\infty$},
  ylabel={pressure coefficient magnitude $|C_p|$},
  xmin=0.4, xmax=0.9, ymin=0, ymax=1.6,
  axis lines=left,
  grid=major, grid style={enggray!20},
  tick label style={font=\scriptsize},
  label style={font=\small},
  legend style={font=\scriptsize, at={(0.03,0.97)}, anchor=north west},
  legend cell align=left,
]
% Corrected suction-peak Cp magnitude: 0.43 / sqrt(1-M^2)
\addplot[engblue, very thick, domain=0.40:0.86, samples=120]
  { 0.43 / sqrt(1 - x*x) };
\addlegendentry{corrected suction peak $|C_{p}|=|C_{p,0}|/\sqrt{1-M_\infty^2}$}
% Critical pressure coefficient magnitude:
% Cp* = (2/(1.4 M^2)) [ ((1+0.2 M^2)/1.2)^3.5 - 1 ], which is negative;
% plot its magnitude = -Cp*.
\addplot[engred, very thick, domain=0.40:0.88, samples=120]
  { -( (2/(1.4*x*x)) * ( ((1 + 0.2*x*x)/1.2)^3.5 - 1 ) ) };
\addlegendentry{critical $|C_p^{*}|$ (local $M=1$)}
% Intersection marker near M_cr = 0.71 (from the worked example)
\fill[engblack] (axis cs:0.71,0.611) circle (1.8pt);
\draw[enggray, dashed] (axis cs:0.71,0) -- (axis cs:0.71,0.611);
\node[engblack, font=\scriptsize, anchor=west] at (axis cs:0.72,0.75)
  {$M_{\mathrm{cr}}\approx 0.71$};
\end{axis}
\end{tikzpicture}
\caption[Graphical construction of the critical Mach number]{The
critical Mach number is where the compressibility-corrected
suction-peak pressure coefficient (blue) first reaches the critical
value $C_p^{*}$ (red), the pressure coefficient at which the local Mach
number equals one. Both are plotted as magnitudes so the crossing reads
directly. The blue curve rises with $M_\infty$ through the
Prandtl-Glauert factor $1/\sqrt{1-M_\infty^2}$; the red curve falls as
$M_\infty$ climbs, because a faster freestream needs a smaller pressure
drop to reach sonic conditions. For the thin section drawn here
($|C_{p,0}|=0.43$ at the suction peak) the two meet at
$M_{\mathrm{cr}}\approx 0.71$. A section with a shallower suction peak
shifts the blue curve down, so it meets the red curve further to the
right, raising the critical Mach number and, with it, the
drag-divergence Mach number. Flattening the upper-surface pressure is
what a supercritical section does to push this crossing toward unity.}
\label{fig:vol03ch13:critical-mach}
\end{figure}
